% !PRINTED 215
\weberpage{215}
%<<CONT
$\varrho$ in $\varrho^m$ übergeht, wenn $m$ der Determinante $(a\partial - bc)$ nach dem Modul $n$ kongruent ist.

4. \gesperrt{Wir wollen schließlich noch die Zahl $\nu$, d.\,h. den Grad der Gruppe $\mathfrak{B}$ bestimmen.}

$\varphi(n)$ hat, wie bekannt, den Ausdruck (Bd.\,I, § 140)

\weq{21}{\varphi(n) = n\varPi\left(1 - \frac{1}{p}\right),}

wenn das Produktzeichen $\varPi$ sich auf alle in $n$ aufgehenden, voneinander verschiedenen Primzahlen $p$ erstreckt.

Die Zahl $\nu$ ist gleich der Anzahl der nach $n$ inkongruenten Zahlensysteme $\alpha, \beta, \gamma, \delta$, die der Bedingung

\weq{22}{\alpha\delta - \beta\gamma \equiv 1 \pmod{n}}

genügen. Wir fragen zunächst nach der Anzahl der Paare $\alpha, \beta$, die mit $n$ keinen gemeinschaftlichen Teiler haben, und bezeichnen diese mit $\chi(n)$.

Ist $n = n'n''$ und $n'$ relativ prim zu $n''$, so kann man aus jeder Kombination eines zu $n'$ gehörigen Zahlenpaares $\alpha', \beta'$ mit einem zu $n''$ gehörigen Zahlenpaar $\alpha'', \beta''$ ein zu $n$ gehöriges

\[
\alpha = n''\alpha' + n'\alpha'', \quad \beta = n''\beta' + n'\beta''
\]

herleiten, und man erhält auf diese Weise alle Zahlenpaare $\alpha, \beta$ und jedes nur einmal. Daraus folgt:

\weq{23}{\chi(n) = \chi(n')\chi(n'').}

Es ist also noch $\chi(p^\pi)$, d.\,h. $\chi(n)$ für eine Primzahlpotenz $n = p^\pi$ zu bestimmen.

Setzen wir zunächst für $\alpha, \beta$ alle modulo $p^\pi$ verschiedenen Zahlen, so ist diese Anzahl $p^{2\pi}$. Hiervon sind aber alle die Paare wegzulassen, bei denen $\alpha$ \gesperrt{und} $\beta$ durch $p$ teilbar sind, deren Anzahl $p^{2\pi-2}$ beträgt, so daß

\[
\chi(p^\pi) = p^{2\pi} - p^{2\pi-2},
\]

und folglich nach (23) allgemein

\weq{24}{\chi(n) = n^2\varPi\left(1 - \frac{1}{p^2}\right)}

folgt, wenn $p$ die in $n$ aufgehenden Primzahlen durchläuft.

Jedes Zahlenpaar $\alpha, \beta$ läßt sich durch Vermehrung um Vielfache von $n$ in ein solches verwandeln, deren Zahlen unter sich relativ prim sind, und dann läßt sich $\gamma, \delta$ so bestimmen, daß

\[
\alpha\delta - \beta\gamma = 1
\]
%CONT>>
